\section{Bessel functions at a real place}\label{sec:bessel}

The Poincar\'e series of \cref{sec:moment} are built from Bessel functions at the infinite places: vector-valued functions on the real
unitary group $G_1$ of $Y_0$ that transform on the left by the generic character of the unipotent radical $U$ of the Bessel subgroup and by an
irreducible representation of $H=\U(X_0)$, and that lie in a discrete series. This section constructs them, with their normalisation, Schur's
bound for their derivatives, their decay, a cut of bounded support, and the \emph{open cell bound} for integrals of a Bessel function against its
translates. The local band limits at one place (\cref{sec:frames}) complete the list of properties that the filling theorem uses
(\cref{def:bessel-system}).

\subsection{The group and its frame}\label{ss:bessel-group}
Let $G_1$ be the complex unitary group of the hermitian form $J_1$ made of $\varrho$ hyperbolic planes $\langle e_a,e^*_a\rangle$, the standard positive form
on $X_0=\C^\kappa$, $\#\kappa=k$, and a vector $z$ with $h(z,z)=-1$; its rank is $r=k+2\varrho+1$. The \emph{frame} is $f=(e_1,\dots,e_\varrho,z)$, of $n=\varrho+1$
vectors. The \emph{Bessel unipotent group} $U$ consists of the isometries that are block upper unitriangular for the levels of the basis; $U^0\subseteq U$ fixes
the frame; $H=\U(X_0)$ is embedded on $X_0$; and the \emph{Bessel subgroup} is $R=UH$, a semidirect product, with the Haar measure $du\,dh$. The
\emph{length} of $y\in G_1$ is $\ell(y)=\inf_{s\in R}(\norm{sy}+1)$, and the \emph{frame box} of radius $\rho$ consists of the $y$ whose \emph{frame vectors}
$y^{-1}f_a$ have norm at most $\rho$.
\lean{New.Shape.G₁, New.Shape.J₁, New.Shape.besselFrame, New.Shape.Form.besselU, New.Shape.Form.besselU0, New.Shape.Form.besselH, New.Shape.Form.besselS, New.Shape.embedH_mem_besselS, New.Shape.embedH_conj_mem_besselU_iff, New.Shape.besselGraph, New.Shape.mem_besselGraph, New.Shape.mem_besselGraph_iff_mul_inv_mem, New.Shape.besselGraphChart, New.Shape.besselGraphMeasure, New.Shape.besselGraphMeasure_finiteOnCompacts, New.Shape.besselGraphMeasure_isMulRightInvariant, New.Shape.besselGraphMeasure_isMulLeftInvariant, New.DiscreteSeries.besselGraph_fst_mem_besselS, New.besselLength, New.besselLength_nonneg, New.frameBox, New.frameVector, New.continuous_frameVector, New.frameVector_mul_of_mem_besselFixer, New.frameVector_eq_sum, New.besselQuotMeasure}

Throughout, the Euclidean product is the one for which $e_a,x_i,z,e^*_a$ is an orthonormal basis, norms of vectors are Euclidean, and norms of matrices are
operator norms; $h(v,w)=\langle v,J_1w\rangle$. Since $J_1$ is a unitary involution, every $g\in G_1$ has $g^{-1}=J_1g^*J_1$, hence $\norm{g^{-1}}=\norm g$ and
$1\le\norm g$. An element $s\in R$ maps $f_b$ to $f_b$ plus a combination of the $f_a$ with $a<b$: its block $m(s)$ on the frame is upper unitriangular,
and elements of $H$ and of $U^0$ fix the frame.

\begin{lemma}[Frame vectors]\label{lem:frame-vectors}
Let $y,\gamma\in G_1$.
\begin{enumerate}[label=\textup{(\alph*)}]
\item $\norm{y^{-1}f_a}\le\norm y$; for $s\in U$, $y^{-1}f_a=\sum_bm(s)_{ba}\,(sy)^{-1}f_b$; for $s$ fixing the frame, $(sy)^{-1}f_a=y^{-1}f_a$.
\item The matrix $Z(\gamma)$ of the products $h(f_a,\gamma f_b)$ satisfies
\[
  Z(\gamma)_{ab}=h\bigl(y^{-1}f_a,\,y^{-1}\gamma f_b\bigr),\qquad\abs{Z(\gamma)_{ab}}\le\norm{y^{-1}f_a}\,\norm{y^{-1}\gamma f_b}.
\] For $s$ fixing the frame,
$Z(s\gamma)=Z(\gamma s)=Z(\gamma)$, and for $s\in U$, $Z(\gamma)=m(s)^*Z(s\gamma)$.
\end{enumerate}
\end{lemma}

\begin{proof}
(a) The frame vectors $f_a$ are unit vectors, and $y^{-1}$ applied to a unit vector has norm at most $\norm{y^{-1}}=\norm y$. Next $y^{-1}=(sy)^{-1}s$ and $sf_a=\sum_bm(s)_{ba}f_b$, so
$y^{-1}f_a=(sy)^{-1}sf_a=\sum_bm(s)_{ba}(sy)^{-1}f_b$; if $s$ fixes the frame, $m(s)=1$.

(b) The isometry $y^{-1}$ preserves $h$, which gives the first identity. For the inequality, $\abs{h(v,w)}=\abs{\langle v,J_1w\rangle}\le\norm v\,\norm{J_1w}=\norm v\,\norm w$, because $J_1$ is unitary. If $s$ fixes the frame, then
$Z(s\gamma)_{ab}=h(s^{-1}f_a,\gamma f_b)=h(f_a,\gamma f_b)$ and $Z(\gamma s)_{ab}=h(f_a,\gamma sf_b)=h(f_a,\gamma f_b)$. For $s\in U$, the frame block is multiplicative on $R$, since $R$ preserves the span of the frame, so
$s^{-1}f_a=\sum_cm(s)^{-1}_{ca}f_c$; as $h$ is conjugate linear in the first variable,
\[
  Z(s\gamma)_{ab}=h(s^{-1}f_a,\gamma f_b)=\sum_c\overline{m(s)^{-1}_{ca}}\,Z(\gamma)_{cb},
\]
that is $Z(s\gamma)=(m(s)^{-1})^*Z(\gamma)$, or $Z(\gamma)=m(s)^*Z(s\gamma)$.
\end{proof}

\begin{proposition}[Length and frames]\label{prop:length-frames}
\begin{enumerate}[label=\textup{(\alph*)}]
\item If the vectors $y^{-1}e_a$, $y^{-1}z$ and $y^{-1}e^*_a$ have norm at most $M\ge1$, then $\norm y+1\le8(\varrho+1)M$. The length is continuous, left invariant under $R$, and
$\ell(y)\le\norm{x^{-1}y}\,\ell(x)$.
\item The Gram--Schmidt vectors of the frame vectors are left invariant under $R$; if $\ell(y)\le M$ their norms lie in $[M^{-1},M]$, and conversely, if they lie in $[M^{-1},M]$
with $M\ge1$, then $\ell(y)\le8(\varrho+1)M$. The triangular change of basis depends continuously and measurably on $y$.
\item The \emph{discriminant} of a $(\varrho+1)\times(\varrho+1)$ matrix $Z$ with leading minors $D_l$ is $\mathfrak g(Z)=\abs{D_{\varrho+1}}^2-\abs{D_\varrho}^2$; up to sign it is the Gram
determinant of the frame and its image under an isometry with frame Gram matrix $Z$.
\end{enumerate}
\end{proposition}
\lean{New.norm_frameVector_le, New.norm_sum_le_mul_norm_sum_frameVector, New.besselLength_le_norm_mul_add_one, New.RealPlace.exists_norm_le_of_besselLength_le, New.besselLength_besselS_mul, New.besselLength_le_norm_inv_mul_mul_besselLength, New.continuous_besselLength, New.norm_J₁, New.norm_sum_rankOne_col, New.norm_add_one_le_of_mem_frameBox, New.norm_inner_J₁_le, New.realFrameGram_eq_inner, New.realFrameGram_mul_of_mem_besselFixer, New.realFrameGram_eq_frameMatrix_mul, New.norm_realFrameGram_le_mul, New.gramSchmidt_eq_of_sub_mem_span, New.orthoFrame, New.gramSchmidtMatrix, New.inner_orthoFrame_eq_zero_of_ne, New.frameVector_eq_sum_orthoFrame, New.exists_frameVector_mul_eq_orthoFrame, New.orthoFrame_besselS_mul, New.norm_orthoFrame_le, New.one_le_norm_mul_norm_orthoFrame, New.orthoFrame_ne_zero, New.norm_orthoFrame_mem_Icc, New.linearIndependent_orthoFrame, New.gramSchmidtMatrix_besselS_mul, New.continuous_orthoFrame, New.continuous_gramSchmidtMatrix, New.measurable_gramSchmidtMatrix, New.besselLength_le_of_norm_orthoFrame, New.besselLength_le_of_dual, New.frameDisc, New.realFrameGram, New.RealPlace.gramDiscReal, New.framePair, New.isotropicFramePair, New.det_eq_add_of_row_eq_add, New.det_fromBlocks_single, New.gram_framePair, New.det_gram_framePair, New.linkMatrix, New.link, New.latticeMiddleGram, New.middleInclusionMatrix, New.fromCols_middleInclusionMatrix, New.fromBlocks_scaleAt_mul_eq_one, New.fromBlocks_one_mul_eq_one, New.besselLatticeBasis, New.gram_besselLatticeBasis}

\begin{proof}
(a) Put $A=y^{-1}$. The Euclidean product is $\langle v,w\rangle=h(v,w)+\langle v,(1-J_1)w\rangle$, and $AJ_1A^*=J_1$. Hence $AA^*=J_1+A(1-J_1)A^*$ and
\[
  \norm A^2=\norm{AA^*}\le1+\norm{A(1-J_1)}\,\norm A .
\]
The columns of $1-J_1$ are $0$ at the $x_i$, $2z$ at $z$, and $\pm(e_a-e^*_a)$ at $e_a$ and $e^*_a$; so the columns of $A(1-J_1)$ have norm at most $2M$, and $\norm{A(1-J_1)}\le(4\varrho+2)M$.
Dividing by $\norm A\ge1$ gives $\norm A\le1+(4\varrho+2)M$, and $\norm y+1=\norm A+1\le2+(4\varrho+2)M\le8(\varrho+1)M$. Left invariance holds by definition. For $s\in R$,
$\norm{sy}\le\norm{sx}\norm{x^{-1}y}$ and $1\le\norm{x^{-1}y}$, so $\norm{sy}+1\le(\norm{sx}+1)\norm{x^{-1}y}$; take the infimum over $s$. Applying this with $x$ and $y$ exchanged,
$\ell(x)/\norm{y^{-1}x}\le\ell(y)\le\norm{x^{-1}y}\ell(x)$, and both bounds tend to $\ell(x)$ as $y\to x$.

(b) Let $g_a(y)$ be the Gram--Schmidt vectors of $y^{-1}f_1,\dots,y^{-1}f_n$ for the Euclidean product (not normalised). If two families $v$, $w$ satisfy $v_a-w_a\in\operatorname{span}(w_b:b<a)$ for all
$a$, they have the same Gram--Schmidt vectors: by induction the earlier ones agree, and the difference of the $a$-th vectors lies in the span $B$ of the earlier ones and is orthogonal to $B$, hence $0$.
For $s\in R$, (a) of \cref{lem:frame-vectors} applied to $s^{-1}$ shows that the frame vectors of $sy$ differ from those of $y$ by combinations of earlier ones, so $g(sy)=g(y)$.
Since $g_a$ is $y^{-1}f_a$ minus its projection to the earlier vectors, $\norm{g_a}\le\norm{y^{-1}f_a}\le\norm y$; and $yg_a$ is $f_a$ plus a combination of $f_b$, $b<a$, so $1\le\norm{yg_a}\le\norm y\norm{g_a}$.
If $\ell(y)\le M$, some $s\in R$ has $\norm{sy}<M$, and these bounds at $sy$, together with $g(sy)=g(y)$, give $\norm{g_a}\in[M^{-1},M]$.

Conversely let $\norm{g_a}\in[M^{-1},M]$. The Gram--Schmidt matrix is upper unitriangular, and is the frame block of an element $u\in U$ with $(uy)^{-1}f_a=g_a$. Put $w_a=J_1g_a/\norm{g_a}^2$ for $a\le\varrho$.
As $h(v,J_1w)=\langle v,w\rangle$, we get $h(g_b,w_a)=\delta_{ab}$; and $h(w_a,w_b)=h(g_a,g_b)/(\norm{g_a}^2\norm{g_b}^2)=0$ for $a,b<n$, since the $g_a$ are the images of the isotropic $e_a$ under an isometry. Hence $(w_a)_{a\le\varrho}$ is an isotropic system dual to the
frame $(g_b)$, and $\norm{w_a}=\norm{g_a}^{-1}\le M$. The images of $(e^*_a)$ under $(uy)^{-1}$ form another such system; two such systems differ by an element of $U^0$ (Witt's theorem with the frame fixed), so
$w_a=(uy)^{-1}u_0e^*_a$ for some $u_0\in U^0$, and $y'=u_0^{-1}uy$ has frame vectors $g_a$ and $y'^{-1}e^*_a=w_a$. By (a), $\ell(y)=\ell(y')\le\norm{y'}+1\le8(\varrho+1)M$. Continuity of $g$ and of the
triangular matrix follows along the Gram--Schmidt recursion, which divides only by the nonzero numbers $\norm{g_a}^2$.

(c) Let $D=\mathrm{diag}(0,\dots,0,-1)$ be the Gram matrix of the frame. For $h$, the Gram matrix of $(f,\gamma f)$ has the diagonal blocks $D$ and the off-diagonal blocks $Z$ and $Z^*$, $Z=Z(\gamma)$. Exchange the two block columns,
at the cost of the sign $(-1)^n$. For square $Z$, $W$ and the matrix $E$ with the single entry $c$ at $(t,t)$, $t$ the last index, expanding along the two rows of index $t$ gives
$\det\begin{pmatrix}Z&E\\E&W\end{pmatrix}=\det Z\det W-c^2\,\mathrm{adj}(Z)_{tt}\,\mathrm{adj}(W)_{tt}$. With $W=Z^*$, $c=-1$ and $\mathrm{adj}(Z)_{tt}=D_\varrho(Z)$ the determinant is $\abs{\det Z}^2-\abs{D_\varrho}^2=\mathfrak g(Z)$.
\end{proof}

For invertible $Z$ let $\xi(Z)=(Z^{-1})_{tt}$, the last diagonal entry of the inverse, and $\nu(Z)^2=1-\abs{\xi(Z)}^2$. By Cramer's rule $\xi=D_\varrho/\det Z$, so
\begin{equation}\label{eq:disc-nu}
  \mathfrak g(Z)=\abs{\det Z}^2\,\nu(Z)^2 .
\end{equation}
The \emph{open cell} $\Omega$ is the set of invertible $Z$ with $\abs{\xi(Z)}<1$; by \eqref{eq:disc-nu} it is the open set $\{\mathfrak g>0\}$.

\subsection{Bessel functions}\label{ss:bessel-functions}
A \emph{Bessel function} of type $\pi$, a unitary representation of $H$ on $\C^I$, is a bounded measurable $W\colon G_1\to\C^I$ with $W(uhy)=\psi(u)\pi(h)W(y)$ for $u\in U$,
$h\in H$, where $\psi(u)=\exp(\sigma(u)-\bar\sigma(u))$ and $\sigma(u)$ is the sum of the entries of $u$ just above the diagonal on the frame. For a doubled
Harish-Chandra parameter $l=2\lambda$ of large gap with discrete series coefficient $c_\lambda$ and formal degree $\deg$, the \emph{Bessel average} at a point $x$ and a
unit vector $v$ is $W_x(y)=\int_Rc_\lambda(x^{-1}sy)\,\pi(s)^{-1}v\,ds$.
\lean{New.IsBesselFunction, New.DiscreteSeries.besselIntegrand, New.DiscreteSeries.besselAverage, New.DiscreteSeries.besselAverage, New.DiscreteSeries.UnitarySignFrame, New.Filling.IsDiscreteSeriesBessel, New.convolve}

Here $\pi$ is extended to $R$ by $\pi(uh)=\psi(u)\pi(h)$; this is a unitary representation of $R$, since $\sigma$ is additive on $U$ and conjugation by $H$, which acts only on $X_0$, does not change the entries
above the diagonal on the frame. The coefficient $c_\lambda$ is written on $G_1$ through a unitary frame in which the form is diagonal (\cref{app:ds}); conjugation by this frame preserves operator norms. We write
$F*k$ for the right convolution $(F*k)(y)=\int_{G_1}k(a^{-1})F(ya)\,da$, and $e=\deg\cdot c_\lambda$. Since $c_\lambda(a^{-1})=\overline{c_\lambda(a)}$, the condition $F*e=F$ says that $F$ is fixed by the idempotent of the discrete
series, and Schur's orthogonality relations give $e*e=e$.

\begin{lemma}[The Bessel subgroup and unfolding]\label{lem:bessel-group}
\begin{enumerate}[label=\textup{(\alph*)}]
\item The measure $du\,dh$ on $R$ is left and right invariant.
\item For measurable $f\ge0$ on $G_1$ and every $x$, $\int_{G_1}f=\int_{U\backslash G_1}\int_Rf(xsy)\,ds\,dy$.
\item If $l$ has gap $500r^3+A$, then $\abs{c_\lambda(g)}\le(2/\norm g)^{j+9r^2}$ for $j\le A$, and
\[
  \int_R\Bigl(\frac2{\norm{x^{-1}sy}}\Bigr)^{9r^2}ds\le16^{9r^2}\,V(x),\qquad V(x)=\vol\{s\in R:\norm{x^{-1}sx-1}\le1\},
\]
uniformly in $y$; and $0<V(x)<\infty$.
\end{enumerate}
\end{lemma}

\begin{proof}
(a) Write the elements of $R$ as $hu$ with $h\in H$, $u\in U$. Right multiplication by $h'u'$ sends $hu$ to $(hh')(h'^{-1}uh'u')$, that is $(h,u)$ to $(hh',h'^{-1}uh'u')$: a right translation of $H$, and on $U$ the conjugation by $h'^{-1}$,
which preserves $du$ because $H$ acts on the entries of $U$ by unitary changes of coordinates, followed by a right translation of $U$. Left multiplication by $h'u'$ sends $hu$ to $(h'h)(h^{-1}u'h\,u)$, a left translation of $H$ and, over
each $h$, a left translation of $U$. Haar measures of $H$ (compact) and of $U$ (unipotent) are left and right invariant, so $du\,dh$ is.

(b) The quotient measure on $U\backslash G_1$ is defined by $\int_{G_1}f=\int_{U\backslash G_1}\int_Uf(uy)\,du\,dy$. Writing $s=hu$ and using this for the function $g\mapsto f(xhg)$,
\begin{align*}
  \int_{U\backslash G_1}\int_Rf(xsy)\,ds\,dy&=\int_H\int_{U\backslash G_1}\int_Uf(xhuy)\,du\,dy\,dh\\
  &=\int_H\int_{G_1}f(xhg)\,dg\,dh=\int_{G_1}f,
\end{align*}
by left invariance of Haar measure on $G_1$ and because $H$ has mass one.

(c) The first bound is \cref{prop:A-decay}(a) with $c=j+9r^2$, which is allowed since $9r^2\le500r^3$ and $j\le A$. For the integral put $J(s)=x^{-1}sx$, a homomorphism from $R$ to $G_1$, viewed in $\U(\varepsilon)$ of size $r$
through the unitary frame, and $y'=x^{-1}y$; then $x^{-1}sy=J(s)y'$. The group $R$ is closed ($U$ is defined by polynomial conditions and $H$ is compact), so $J$ is proper, and $ds$ is right invariant by (a). The last statement of
\cref{lem:A-covering} with $d=r$ and $c=9r^2>8r^2$ gives $\int_R(2/\norm{J(s)y'})^{9r^2}ds\le16^{9r^2}\,\vol\{s:\norm{J(s)-1}\le1\}$, uniformly in $y'$. The set in $V(x)$ is the preimage of a compact set under the
proper map $J$, hence compact, so $V(x)<\infty$; it contains the open neighbourhood $\{\norm{J(s)-1}<1\}$ of $1$, so $V(x)>0$.
\end{proof}

\begin{theorem}[Bessel averages]\label{thm:bessel-average}
Let $l$ have gap at least $500r^3$ plus an excess, and let $\pi$ be continuous.
\begin{enumerate}[label=\textup{(\alph*)}]
\item The integrand is integrable and $W_x$ is a continuous Bessel function of type $\pi$, reproduced by the idempotent $\deg\cdot c_\lambda$; and
$\deg\norm{W_x}^2_{L^2(U\backslash G_1)}=\langle v,W_x(x)\rangle$.
\item If $\langle v,\pi(h)v\rangle$ is the coefficient of a parameter $u$ of $H$ and the volume of $\{s\in R:\norm{x^{-1}sx-1}\le1\}$ is at most $P$ times the modulus of the period
$\int_Rc_\lambda(x^{-1}sx)\langle v,\pi(s)^{-1}v\rangle ds$, then it is at most $P\deg\norm{W_x}^2$.
\item (Stationary point.) If $l$ is \emph{steep} (lowest $K$-type of weight $0$ on $X$ and decreasing weight in $[-3B,-B]$ on the other coordinates, with alternating signs) and
the parameters are bounded with excess at least $10^{30}((r+k)^2\log\beta)^3$, there is $x$ with $\ell(x)\le b^{5n}$ at which this volume is at most
$\beta^{20(r+k)^2}$ times the modulus of the period.
\end{enumerate}
\end{theorem}
\lean{New.DiscreteSeries.integrable_besselIntegrand, New.DiscreteSeries.continuous_besselIntegrand, New.DiscreteSeries.enorm_besselIntegrand, New.DiscreteSeries.lintegral_enorm_besselIntegrand_le, New.DiscreteSeries.besselAverage_mul, New.DiscreteSeries.isBesselFunction_besselAverage, New.DiscreteSeries.convolve_besselAverage, New.DiscreteSeries.enorm_besselAverage_le, New.DiscreteSeries.integrable_besselAverage, New.DiscreteSeries.formalDegree_mul_integral_norm_sq_besselAverage, New.DiscreteSeries.besselBallVolume, New.DiscreteSeries.besselBallVolume_ne_top, New.DiscreteSeries.besselBallVolume_ne_zero, New.DiscreteSeries.lintegral_decay_le_besselBallVolume, New.DiscreteSeries.norm_dsCoeffInFrame_le, New.DiscreteSeries.besselBallVolume_le_of_period, New.DiscreteSeries.IsSteep.exists_stationary_point, New.DiscreteSeries.lintegral_unfold_besselGraphMeasure, New.DiscreteSeries.measurePreserving_unfold_besselGraphMeasure, New.DiscreteSeries.exists_sub_one_mul_frameCols_eq, New.DiscreteSeries.exists_sub_one_mul_definiteCols_eq, New.DiscreteSeries.exists_besselGraph_fst_eq, New.DiscreteSeries.besselPairHom, New.DiscreteSeries.isProperMap_besselPairHom_fst, New.DiscreteSeries.subOnePair_besselPairHom_fst, New.DiscreteSeries.besselPairHom_snd_val, New.DiscreteSeries.subOnePair_besselPairHom_snd, New.DiscreteSeries.norm_frame_mulVec_le, New.DiscreteSeries.IsPosition.exists_mul_eq_frameCols_and_besselLength_le, New.DiscreteSeries.initCols_eq_of_frame, New.DiscreteSeries.lowersFlag_besselPairHom, New.DiscreteSeries.inducesOn_besselPairHom, New.DiscreteSeries.exists_besselPairHom_fst_eq, New.DiscreteSeries.mul_sqrt_mul_pow_le_pow}

\begin{proof}
(a) As $\pi$ is unitary and $\norm v=1$, the integrand has modulus $\abs{c_\lambda(x^{-1}sy)}\le(2/\norm{x^{-1}sy})^{9r^2}$, which by \cref{lem:bessel-group}(c) is integrable over $R$ with integral at most $16^{9r^2}V(x)$,
uniformly in $y$. Hence $W_x$ is bounded; it is continuous because the integrand is jointly continuous and dominated. For $s\in R$, the substitution $t\mapsto ts^{-1}$, allowed by right invariance, gives
\[
  W_x(sy)=\int_Rc_\lambda(x^{-1}t\,sy)\,\pi(t)^{-1}v\,dt=\int_Rc_\lambda(x^{-1}t'y)\,\pi(t's^{-1})^{-1}v\,dt'=\pi(s)W_x(y),
\]
so $W_x$ is a Bessel function of type $\pi$. For the reproduction, exchange the integrals in $(W_x*e)(y)=\int_{G_1}e(a^{-1})W_x(ya)\,da$; the double integral converges absolutely by the covering bound and the
integrability of $c_\lambda$ over $G_1$ (\cref{lem:A-covering}). The inner integral is $\int_{G_1}c_\lambda(x^{-1}sya)\,e(a^{-1})\,da=(c_\lambda*e)(x^{-1}sy)=c_\lambda(x^{-1}sy)$.

For the norm, insert the definition of $W_x(y)$ in one factor of $\langle W_x(y),W_x(y)\rangle$ and move $\pi(s)^{-1}=\pi(s)^*$ to the other side:
\[
  \langle W_x(y),W_x(y)\rangle=\int_R\overline{c_\lambda(x^{-1}sy)}\,\langle v,\pi(s)W_x(y)\rangle\,ds=\int_R\Phi(sy)\,ds,
\]
where $\Phi(g)=\overline{c_\lambda(x^{-1}g)}\,\langle v,W_x(g)\rangle$.
By \cref{lem:bessel-group}(b) with $x=1$, $\norm{W_x}^2=\int_{G_1}\Phi$. Moreover $\overline{c_\lambda(x^{-1}g)}=c_\lambda(g^{-1}x)$, and with $g=xa$, $\int_{G_1}c_\lambda(g^{-1}x)W_x(g)\,dg=\deg^{-1}(W_x*e)(x)=\deg^{-1}W_x(x)$. Hence
$\deg\norm{W_x}^2=\langle v,W_x(x)\rangle$.

(b) By definition $\langle v,W_x(x)\rangle=\int_Rc_\lambda(x^{-1}sx)\langle v,\pi(s)^{-1}v\rangle\,ds$ is the period, and by (a) it equals $\deg\norm{W_x}^2\ge0$. The hypothesis on $\langle v,\pi(h)v\rangle$ identifies the integrand with the Laplace
integrand of \cref{thm:A-period}: for $s=uh$, $\langle v,\pi(s)^{-1}v\rangle=\overline{\psi(u)}\,\overline{c_u(h)}=\exp(\bar\sigma-\sigma)\,\overline{c_u(h)}$.

(c) Let $\mu$ be the lowest weight of $l$ and $\mu'$ that of $u$, which is $a$ at one coordinate and $0$ on $X$. Apply \cref{prop:A-stationary} with $d=\varrho$, $L=2n$ coordinates of $G_1$ outside $X$ and $L'=1$ coordinate
of $H$ outside $X$; then $L=L'+2d+1$ and there is a stationary position $(v_1,\dots,v_n;i_0)$ with $v_n$ of norm $-1$, whose entries, and those of $\mu^av_1$ for $a\le2\varrho$, are at most $M=(3B+a)^{2n}$. The vectors
$w_a=\mu^av_1$, $a\le2\varrho$, continue $v_1,\dots,v_n$, and by the residue computation in the proof of \cref{prop:A-stationary} their products are $0$ for $a+b<2\varrho$ and $-1$ for $a+b=2\varrho$. The columns of $i_0$ and
$v_1,\dots,v_n$ have the Gram matrix of the basis of $X_0$ and of the frame, and both families are independent; by Witt's theorem there is $x\in G_1$ carrying $i_0$ to the inclusion of $X_0$ and $v$ to the frame
(after the unitary frame of signs). The vectors $x^{-1}$-images of the $e^*_a$ are then given by the $w_a$, of norm at most $\sqrt rM$, and the argument of \cref{prop:length-frames}(b) gives $\ell(x)\le8n\sqrt rM\le b^{5n}$,
using $3B+a\le b^2$ and $8nr\le b$. For this $x$, conjugation $s\mapsto x^{-1}sx$ maps $R$ properly onto the group of the position: $x^{-1}sx-1$ lowers the flag of the $v_j$, with sum of coefficients
$\sigma(s)$, and induces on $i_0$ the image of $s$ in $H$; conversely every element of the group of the position arises. \Cref{thm:A-period} with $N=r+k$ bounds the volume of the $s$ with $\norm{x^{-1}sx}\le\beta+1$, which
contains the unit ball since $H$ is compact, by $\beta^{20(r+k)^2}$ times the modulus of the period.
\end{proof}

\begin{theorem}[Schur's bound, decay and cut]\label{thm:bessel-schur}
Let $F$ be a Bessel function of type $\pi$, integrable over $U\backslash G_1$ and reproduced by $\deg\cdot c_\lambda$, for $l$ of gap $500r^3+A$ with entries bounded by $L$.
\begin{enumerate}[label=\textup{(\alph*)}]
\item (Schur.) $F$ has right derivatives $V_w$ along all words $w$ in the matrix units, again Bessel functions of type $\pi$; for $\abs w\le2n^2$ and integrable $\varphi$,
$\norm{g\mapsto\int\varphi(y)V_w(gy)dy}^2_{L^2(U\backslash G_1)}\le B^{2\abs w}\norm F_2^2\norm\varphi_2^2/\deg$ whenever $B\ge(r+1)^4(L+r+1+2n^2)$.
\item (Decay.) Let $x$ have $\ell(x)\le b^{5n}$ and $b^m\deg\norm{W_x}^2\ge$ the ball volume, with $b\ge2000$, $r+L+1\le b^4$, $m\le301r^2$, $3r^3f\le A$. The normalised
$F=cW_x$ with $\norm F_2^2=\#I$ and its cut $F^\flat=\chi(\ell/b^{6n+2})F$, $\chi$ a fixed cutoff, have a \emph{decay majorant}: a left $R$-invariant $w\le a<\infty$ with
$\max(1,\ell/b^{6n})^{3r^3f}\abs{F}$, $\max(1,\ell/b^{6n})^{3r^3f}\abs{F^\flat}\le w$ and $a\int w\le b^{400r^2}\#I$.
\item (Cut pair.) Then $F$ and $F^\flat$ are Bessel functions, $\norm F_2^2=\#I$, $\norm{F^\flat}_2^2\le\#I$, $F^\flat$ vanishes where $\ell>b^{6n+2}$, $F^\flat=F$ where $\ell\le\frac12b^{6n+2}$, and
$\norm{F-F^\flat}_2^2\le b^{-2r^3f}\#I$.
\item (Orthonormality.) If $\pi$ is irreducible, the operator $\int_{U\backslash G_1}F\otimes F^*$ is $1$ for $F$ and a scalar in $[0,1]$ for $F^\flat$; so the coordinates of
$F$ are orthonormal in $L^2(U\backslash G_1)$.
\end{enumerate}
\end{theorem}
\lean{New.IsBesselFunction.hasSchurBound_of_convolve_eq, New.isBoundedMeasurable_formalDegree_mul_dsCoeffInFrame, New.convolve_formalDegree_mul_dsCoeffInFrame_self, New.convolve_formalDegree_mul_dsCoeffInFrame_comp_frameEquiv, New.HCParam.HasLargeGap.norm_sq_eq_formalDegree_mul_integral, New.exists_derivative_kernels, New.IsBesselFunction.lintegral_enorm_sq_integral_inner_translate, New.IsBesselFunction.convolve_of_integrable, New.SchurLe, New.RealPlace.HasSchurBound, New.exists_isDecayMajorant_besselAverage, New.exists_isCutPair_besselAverage, New.IsDecayMajorant, New.IsCutPair, New.InCut, New.scaledBesselLength, New.one_le_scaledBesselLength, New.le_scaledBesselLength, New.besselCutoff, New.besselCutoff_mem_Icc, New.besselCutoff_eq_one, New.besselCutoff_eq_zero, New.besselCut, New.besselCut_apply, New.enorm_besselCut_le, New.enorm_besselCut_sub_le, New.inCut_of_besselCut_ne_zero, New.lt_besselLength_of_besselCut_ne, New.IsBesselFunction.besselCut, New.IsDecayMajorant.lintegral_sq_besselCut_sub_le, New.lintegral_besselGraph_mul_eq, New.besselLength_div_le_half_norm, New.scaledBesselLength_pow_mul_enorm_besselAverage_le, New.IsBesselFunction.smul_of_invariant, New.exists_enorm_sq_mul_eq, New.BesselSystem.orthonormal, New.IsBesselFunction.integral_mul_conj_embedH_mul, New.IsBesselFunction.integral_mul_conj_eq, New.IsCutPair.integrable_cut, New.IsCutPair.exists_integral_cut_mul_conj_eq, New.IsCutPair.lintegral_sq_inner_cut_le}

The proof uses two lemmas on the discrete series at one real place: the derivatives of the coefficient along words in the matrix units, and the decay of the coefficient measured against the formal degree.
The coefficient $c_\lambda$ is a smooth vector, and $U(\mathfrak g)c_\lambda$ is the space of its right derivatives, on which a matrix $X$ acts by the derivative along $X$. We norm it by
$\norm x=\deg^{1/2}\norm x_{L^2(\U(\varepsilon))}$, so that $\norm{c_\lambda}=1$, the formal degree being the inverse of $\int\abs{c_\lambda}^2$. Let $M_0$ be the least subspace that contains $c_\lambda$ and is stable under the
matrix units $E_{ab}$ with $\varepsilon_a=\varepsilon_b$, and $M_{s+1}=M_s+\sum_{a,b}E_{ab}M_s$.

\begin{lemma}[Steps of the filtration]\label{lem:filtration-steps}
Let $l$ have gap $500r^3+A$ and entries bounded by $L$.
\begin{enumerate}[label=\textup{(\alph*)}]
\item For $x\in M_s$ and every matrix $X$, $Xx\in M_{s+1}$ and $\norm{Xx}\le(r+1)^4(L+r+1+s)\norm X\norm x$.
\item Let $\beta$ be a bump function on $\U(\varepsilon)$ with $c_\lambda*\beta=c_\lambda$, and $w$ a word of length $s$ in the matrix units, conjugated by the unitary frame. Then $c_\lambda*\partial_w\beta$ lies in $M_s$ and
\[
  \norm{c_\lambda*\partial_w\beta}\le\bigl((r+1)^4(L+r+1+s)\bigr)^s.
\]
\end{enumerate}
\end{lemma}
\lean{New.norm_frame_conj_single, New.exists_mem_filtration_smear_word}

\begin{proof}
(a) $Xx=\sum_{a,b}X_{ab}E_{ab}x$ lies in $M_{s+1}$ by definition. For the norm we bound $\Sigma(y)=\sum_{a,b}\norm{E_{ab}y}^2$ for $y\in M_s$. The $L^2$ product is invariant, so $E_{ab}$ is adjoint to $\varepsilon_a\varepsilon_bE_{ba}$, and
for the Casimir operator $\Omega=\sum_{a,b}E_{ab}E_{ba}$ and its compact part $\Omega_K$, the sum over the pairs with $\varepsilon_a=\varepsilon_b$,
\[
  \langle y,\Omega y\rangle=\sum_{\varepsilon_a=\varepsilon_b}\norm{E_{ba}y}^2-\sum_{\varepsilon_a\ne\varepsilon_b}\norm{E_{ba}y}^2,\qquad \Sigma(y)=2\langle y,\Omega_Ky\rangle-\langle y,\Omega y\rangle.
\]
On the discrete series $\langle y,\Omega y\rangle\ge\kappa_\lambda\norm y^2$ with $\kappa_\lambda=\sum_a(l_a^2-(2\rho)_a^2)/4$, and $\abs{(2\rho)_a}\le r$, so $-\langle y,\Omega y\rangle\le-\kappa_\lambda\norm y^2\le\frac{r^3}4\norm y^2$. On the other side, the generator $c_\lambda$ is the highest
vector of the lowest $K$-type, of weight $\Lambda$ with $\abs{\Lambda_a}\le L+r$. A $K$-type of $M_s$ has a highest vector not orthogonal to some $E_wc_\lambda$ with $w$ a word of length at most $s$, so its highest weight $\mu$ is the
weight of $E_wc_\lambda$, which has entries of modulus at most $N=L+r+s$; on that $K$-type $\Omega_K$ acts by $c_K(\mu)=\sum_a\mu_a(\mu_a+(2\rho_c)_a)\le rN(N+r)$. Hence $\Sigma(y)\le C\norm y^2$ with
$C=2rN(N+r)+r^3/4$. By Cauchy--Schwarz over the $r^2$ pairs and $\abs{X_{ab}}\le\norm X$,
\[
  \norm{Xx}\le\norm X\sum_{a,b}\norm{E_{ab}x}\le\norm X\bigl(r^2C\bigr)^{1/2}\norm x\le(r+1)^4(N+1)\norm X\norm x,
\]
since $r^2C\le(r+1)^8(N+1)^2$, an inequality between polynomials with nonnegative coefficients.

(b) Induction on $w$. For the empty word, $c_\lambda*\beta=c_\lambda\in M_0$ has norm $1$. Let $w=qw'$ with a first letter $q=(a,b)$, and let $M$ be the matrix unit $E_{ab}$ conjugated by the unitary frame; it has
operator norm $1$, as conjugation by a unitary matrix keeps the norm. Convolution on the left commutes with right derivatives, so $c_\lambda*\partial_w\beta$ is the derivative along $M$ of $x=c_\lambda*\partial_{w'}\beta$, that is $Mx$. By induction
$x\in M_{s-1}$ with $\norm x\le((r+1)^4(L+r+s))^{s-1}$, and (a) gives $Mx\in M_s$ with
\[
  \norm{Mx}\le(r+1)^4(L+r+s)\,\norm x\le\bigl((r+1)^4(L+r+1+s)\bigr)^s.\qedhere
\]
\end{proof}

\begin{lemma}[Formal degree and decay]\label{lem:degree-decay}
Let $l$ have gap $500r^3+A$ and entries bounded by $L$, and $K=2000(r+L+1)^6$.
\begin{enumerate}[label=\textup{(\alph*)}]
\item For every $R\ge1$, $\deg\cdot\vol\{g\in G_1:\norm g\le R\}\le4(KR^2)^{2r^2}$.
\item For every integer $M>8r^2$ the function $(2/\norm g)^M$ is integrable on $G_1$, and $\deg\int_{G_1}(2/\norm g)^Mdg\le8K^{2r^2}2^M$. In particular
\begin{equation}\label{eq:degree-decay}
  \deg\int_{G_1}\Bigl(\frac2{\norm g}\Bigr)^Mdg\le b^{50r^2+1}2^M\quad\text{if }b\ge2000,\ r+L+1\le b^4,
\end{equation}
and then, for $M=9r^2$, $\deg\cdot16^M\int_{G_1}(2/\norm g)^Mdg\le b^{99r^2}$.
\item If $b\ge2000$, $r+L+1\le b^4$ and $T\ge1$, then $\deg\cdot\vol\{g:\norm g\le T\}\le T^{9r^2}b^{99r^2}$.
\end{enumerate}
\end{lemma}
\lean{New.norm_frameEquiv, New.lintegral_formalDegree_mul_decay_le, New.formalDegree_mul_lintegral_decay_le, New.formalDegree_mul_volume_ball_le}

\begin{proof}
Conjugation by the unitary frame carries $G_1$ onto $\U(\varepsilon)$, keeps operator norms and carries Haar measure to Haar measure; so we may work on $\U(\varepsilon)$, whose matrices have size $r$. Put $\beta=r+L+1$.

(a) By \cref{prop:A-decay}(c), $\abs{c_\lambda}\ge1/2$ on the set of $g$ with $\norm{g-1}\le(200\beta^5)^{-1}$, so this set has volume at most $\int4\abs{c_\lambda}^2=4/\deg$. \Cref{lem:A-covering}, applied to the identity of $\U(\varepsilon)$ with
$y=1$, $D=200\beta^5$ and $\kappa=KR^2\ge600\beta^5R^2=3DR^2$, gives $\vol\{\norm g\le R\}\le(KR^2)^{2r^2}\cdot4/\deg$.

(b) Since $\norm g\ge1$, the layer cake formula and (a) with $R=t^{-1/M}\ge1$ give
\begin{align*}
  \int\norm g^{-M}dg&=\int_0^1\vol\{g:\norm g^{-M}>t\}\,dt\\
  &\le\frac{4K^{2r^2}}\deg\int_0^1t^{-4r^2/M}dt=\frac{4K^{2r^2}}\deg\cdot\frac M{M-4r^2}.
\end{align*}
As $M>8r^2$, the last factor is less than $2$; this proves the integrability, and multiplication by $\deg\cdot2^M$ gives the bound. If $b\ge2000$ and $r+L+1\le b^4$, then $K\le2000b^{24}\le b^{25}$ and $8\le b$, so
$8K^{2r^2}\le b^{50r^2+1}$, which is \eqref{eq:degree-decay}. For $M=9r^2$, $2^M16^M\le b^{2M}$, and $50r^2+1+18r^2\le99r^2$.

(c) Let $M=9r^2$. On the ball $\norm g\le T$, $1\le(T/\norm g)^M=(T/2)^M(2/\norm g)^M$. Integrating over the ball and using \eqref{eq:degree-decay},
\[
  \deg\cdot\vol\{\norm g\le T\}\le(T/2)^M\deg\int_{G_1}\Bigl(\frac2{\norm g}\Bigr)^Mdg\le T^Mb^{50r^2+1}\le T^{9r^2}b^{99r^2}.\qedhere
\]
\end{proof}

\begin{proof}[Proof of \cref{thm:bessel-schur}]
(a) Let $X_\emptyset=e$ and $X_w=e*\partial_w\beta$ for a word $w$, where $\partial_w$ is the derivative along the letters of $w$ conjugated by the frame. By \cref{lem:filtration-steps}(b), as $\int\abs{X_w}^2/\deg=\norm{c_\lambda*\partial_w\beta}^2$,
$\int\abs{X_w}^2/\deg\le((r+1)^4(L+r+1+\abs w))^{2\abs w}\le B^{2\abs w}$ for $\abs w\le2n^2$, and $X_\emptyset*X_w=X_w$. Put $V_w=F*X_w$; since $F=F*e$, differentiating under the integral shows that along a curve
$\gamma$ with $\gamma'=\gamma Z$ the derivative of $V_w\circ\gamma$ is $\sum_{a,b}Z_{ab}V_{(a,b)w}\circ\gamma$, so the $V_w$ are the right derivatives of $F$. Convolution on the right commutes with left translations, so
each $V_w$ is again a Bessel function of type $\pi$. For the bound, the right translates of $F$ span a space on which $G_1$ acts by a multiple of the discrete series of $l$, because $F=F*e$; on it the operator of an
integrable kernel $k$ has norm at most $\norm k_2/\deg^{1/2}$ (Plancherel's inequality for the discrete series, from Schur's orthogonality). The function $g\mapsto\int\varphi(y)V_w(gy)\,dy$ is the image of $F$ under the
operators of $X_w$ and of $\varphi$; so its squared norm is at most $\norm F_2^2\cdot\deg^{-1}\int\abs{X_w}^2\cdot\deg^{-1}\int\abs\varphi^2\le B^{2\abs w}\norm F_2^2\norm\varphi_2^2/\deg$.

(b) Put $M=9r^2$, $\varphi(g)=(2/\norm g)^M$ and $\Phi(y)=\int_R\varphi(x^{-1}sy)\,ds$, which is left invariant under $R$ by \cref{lem:bessel-group}(a), and write $\tilde\ell(y)=\max(1,\ell(y)/b^{6n})$. Since
$\ell(y)=\ell(sy)\le\norm{x^{-1}sy}\,\ell(x)\le\norm{x^{-1}sy}\,b^{5n}\le\norm{x^{-1}sy}\,b^{6n}/2$ by \cref{prop:length-frames}(a) and $b\ge2$, every $s$ has $\tilde\ell(y)\le\max(1,\norm{x^{-1}sy}/2)$. Let $i=3r^3f\le A$ and $g=x^{-1}sy$. If $\tilde\ell(y)=1$, the decay of \cref{lem:bessel-group}(c) with $j=0$ gives $\tilde\ell(y)^i\abs{c_\lambda(g)}\le\varphi(g)$. Otherwise $\norm g\ge2\tilde\ell(y)>2$, and the decay with $j=i$
gives
\[
  \tilde\ell(y)^i\abs{c_\lambda(g)}\le\Bigl(\frac{\norm g}2\Bigr)^i\Bigl(\frac2{\norm g}\Bigr)^{i+M}=\varphi(g).
\] Integrating over $R$, $\tilde\ell^i\abs{W_x}\le\Phi$. Take $w=\abs c\,\Phi$; as $\abs{F^\flat}\le\abs F$, both bounds of the decay majorant hold. By \cref{lem:bessel-group}(c), $\Phi\le16^MV(x)$, so $w\le a=\abs c\,16^MV(x)$.
By \cref{lem:bessel-group}(b), $\int_{U\backslash G_1}\Phi=\int_{G_1}\varphi$. Using $c^2\norm{W_x}^2=\#I$ and $V(x)\le b^m\deg\norm{W_x}^2$,
\[
  a\int w=c^2\,16^M\,V(x)\int_{G_1}\varphi\le b^m\,\#I\cdot\deg\,16^M\int_{G_1}\varphi\le b^{m+99r^2}\,\#I\le b^{400r^2}\,\#I,
\]
where \cref{lem:degree-decay}(b) gives $\deg\,16^M\int\varphi\le b^{99r^2}$, and $m\le301r^2$.

(c) The constant $c$ exists: $V(x)>0$ and $V(x)\le b^m\deg\norm{W_x}^2$ give $\norm{W_x}_2>0$. The cut is $F^\flat(y)=\chi(\ell(y))F(y)$ with $\chi(t)=\max(0,\min(1,2-2t/R))$, $R=b^{6n+2}$; it is continuous, a Bessel function because $\ell$ is
left invariant under $R$ (\cref{prop:length-frames}(a)), equal to $F$ where $\ell\le R/2$ and zero where $\ell\ge R$, and $\abs{F^\flat}\le\abs F$, so $\norm{F^\flat}_2^2\le\#I$. Where $F^\flat\ne F$ we have $\ell>R/2$, so $\tilde\ell\ge b^2/2\ge b$ and
$b^{3r^3f}\abs{F^\flat-F}\le\tilde\ell^{3r^3f}\abs F\le w\le a$. Multiplying the two bounds $\abs{F^\flat-F}\le b^{-3r^3f}a$ and $\abs{F^\flat-F}\le b^{-3r^3f}w$ and integrating,
\[
  \norm{F^\flat-F}_2^2\le b^{-6r^3f}a\int w\le b^{400r^2-6r^3f}\,\#I\le b^{-2r^3f}\,\#I,
\]
since $100r^2\le r^3\le r^3f$.

(d) For $s\in H$, left translation by $s$ normalises $U$ and preserves its Haar measure, and the characters cancel in $W(y)\otimes W(y)^*$; so $\int_{U\backslash G_1}W(sy)\otimes W(sy)^*=\int W\otimes W^*$. With $W(sy)=\pi(s)W(y)$ this
says that the operator $T=\int W\otimes W^*$ commutes with $\pi(H)$, hence is a scalar by irreducibility, and its trace is $\norm W_2^2$. For $F$ this is $\#I$, so $T=1$; for $F^\flat$ the scalar lies in $[0,\norm{F^\flat}_2^2/\#I]\subseteq[0,1]$.
\end{proof}

\subsection{The open cell}\label{ss:open-cell}
For $\gamma\in G_1$ let $Z(\gamma)$ be the matrix of products of $f$ with $\gamma f$. The \emph{open cell} $\Omega$ is the set of invertible $Z$ whose inverse has last diagonal entry $\xi$
with $\abs\xi<1$; on it $\nu(Z)^2=1-\abs\xi^2$. There is an explicit real-analytic section $Z\mapsto\gamma_Z\in G_1$ with $Z(\gamma_Z)=Z$, with polynomial bounds for it and its
derivative where $Z$, $Z^{-1}$ and $\nu^{-1}$ are bounded (\emph{bounded in the cell}); by Witt's theorem $Z(\gamma)=Z\in\Omega$ determines $\gamma$ up to the fixer $U^0H$ of the frame on
both sides.
\lean{New.RealPlace.invCorner, New.Shape.openCell, New.RealPlace.oneSubNormSqInvCorner, New.RealPlace.IsBoundedInCell, New.RealPlace.diagLast, New.RealPlace.sectionCoeff, New.RealPlace.sectionFrame, New.RealPlace.sectionMatrix, New.RealPlace.mem_openCell, New.RealPlace.invCorner_conjTranspose, New.RealPlace.conjTranspose_mem_openCell_iff, New.RealPlace.gramDiscReal_eq_norm_det_sq_mul, New.RealPlace.mem_openCell_iff, New.RealPlace.continuous_gramDiscReal, New.RealPlace.isOpen_openCell, New.RealPlace.oneSubNormSqInvCorner_pos, New.RealPlace.differentiableAt_sectionMatrix, New.RealPlace.norm_sectionMatrix_le_and_norm_fderiv_le, New.RealPlace.IsBoundedInCell.norm_sectionMatrix_le, New.RealPlace.sectionMatrix_isometry, New.RealPlace.gramSection, New.RealPlace.coe_gramSection, New.RealPlace.measurable_gramSection, New.RealPlace.realFrameGram_gramSection, New.RealPlace.exists_eq_mul_gramSection_mul}

We say that $Z$ is \emph{bounded in the cell by $P$} if $Z\in\Omega$ and $\norm Z$, $\norm{Z^{-1}}$ and $\nu(Z)^{-1}$ are at most $P$. Let $F$ be the matrix with the columns $f$, and $F'$ that with the columns
$e^*_1,\dots,e^*_\varrho,x_{k_0}$ for a fixed index $k_0$. For $Z\in\Omega$ let $\zeta$ be the last row of $Z^{-1}$ and $Y(Z)$ the matrix with last row $(0,\dots,0,-1)$ and other rows those of $-\zeta^*\zeta$, halved outside the last
column. Then $F^\vee(Z)=F\,Y(Z)+F'\,\mathrm{diag}(1,\dots,1,\nu(Z))$ is a family dual to the frame ($F^*J_1F^\vee=1$) with Gram matrix $-\zeta^*\zeta$, the congruence of the Gram matrix $D$ of the frame by $Z^{-1}$. The section $\gamma_Z$ is the isometry
that maps $F$ to $F^\vee(Z)Z$, maps $F^\vee(Z^*)$ to $F(Z^*)^{-1}$ and fixes the $x_i$, $i\ne k_0$; its matrix is a polynomial in the entries of $Z$, $Z^{-1}$, $\nu$ and $\nu^{-1}$, written out in the archive.

\begin{proposition}[The section of the frame Gram map]\label{prop:gram-section}
\begin{enumerate}[label=\textup{(\alph*)}]
\item If $Z\in\Omega$, then $\gamma_Z\in G_1$ and $Z(\gamma_Z)=Z$. The map $Z\mapsto\gamma_Z$ is measurable on all matrices (the identity off $\Omega$) and real differentiable on $\Omega$.
\item If $Z(\gamma)=Z\in\Omega$, then $\gamma=s\gamma_Zs'^{-1}$ with $s$, $s'$ in the fixer $U^0H$ of the frame.
\item If $Z$ is bounded in the cell by $T\ge4$, then $\norm{\gamma_Z}\le T^{18}$ and the real derivative of $Y\mapsto\gamma_Y$ at $Z$ has norm at most $T^{25}$.
\item (Perturbation.) Let $Z$ be bounded in the cell by $T\ge4$ and $\norm{Y-Z}<T^{-6}$. Then $Y$ is invertible with $\norm{Y^{-1}}\le2T$ and $\norm{Y^{-1}-Z^{-1}}\le2T^{-4}$, $\nu(Y)^2\ge(1-8T^{-2})\nu(Z)^2$, and $Y$ is bounded in the cell by
$2T$; if moreover $T\ge8$ and $T\ge n$, then $\mathfrak g(Y)\ge\mathfrak g(Z)/2$.
\item (Sizes at and near a point.) If $n\le M$, the entries of $Z$ are at most $M^2$ and $\mathfrak g(Z)\ge M^{-1}$, then $Z$ is bounded in the cell by every $T\ge M^{8n}$. If moreover $M\ge8n$ and $P\ge M^{9n}$, every $Y$ with
$\norm{Y-Z}<P^{-6}$ has $\mathfrak g(Y)\ge\mathfrak g(Z)/2$ and is bounded in the cell by $P$.
\end{enumerate}
\end{proposition}

\begin{proof}
(a) The products of $F$ with $\gamma_ZF=F^\vee(Z)Z$ are $F^*J_1F^\vee(Z)Z=Z$; the products of $F^\vee$ with $F$ are the identity off the last index, because $F^*J_1F'=\mathrm{diag}(1,\dots,1,0)$ and $F^*J_1F=\mathrm{diag}(0,\dots,0,-1)$, and the factor
$\mathrm{diag}(0,\dots,0,-1)$ keeps only the last row $(0,\dots,0,-1)$ of $Y(Z)$. That $\gamma_Z$ is an isometry is checked in the span of the orthonormal columns $(F\,|\,F')$, where $\gamma_Z$ acts as $yx^{-1}$ for the
matrices $x$, $y$ whose columns are $(FZ^{-1}\,|\,F^\vee(Z^*))$ and $(F^\vee(Z)\,|\,F(Z^*)^{-1})$: these two families have the same Gram matrix, which reduces to polynomial identities in the entries of $Z^{-1}$ given
$\nu^2=1-\abs\xi^2$; on the orthogonal complement $\gamma_Z$ is the identity. The entries are holomorphic in $(Z,W)$ at $W=Z^*$ near $\Omega\times\Omega$, the square root being taken of $1-\xi(Z)\xi(W)$, which lies in the slit plane; restricting to
$W=Z^*$ gives real differentiability, and measurability follows. (b) The families $(F\,|\,\gamma F)$ and $(F\,|\,\gamma_ZF)$ have the same Gram matrix, which is invertible because its determinant is $\pm\mathfrak g(Z)\ne0$
(\cref{prop:length-frames}(c)). By Witt's theorem some $g$ carries one to the other; $g$ fixes the frame, and so does $\gamma_Z^{-1}g\gamma$. (c) Write the matrix of $\gamma_Z$ as the value at $W=Z^*$ of a matrix in two variables $(Z,W)$. If $P\ge8$, $Z$ is bounded in the cell by $P$ and $W\in\Omega$ has $\norm{W^{-1}}\le P$, the formula gives the norm bound $P^{12}$:
$\norm{Y}\le2+8P^2$ for the coefficient matrix, $\abs s\le2$ and $\abs{s^{-1}}\le2P$ for the root $s$ of $1-\xi(Z)\xi(W)$, and the factors $F$, $F'$ have norm one. By (d) every $(Y,W)$ in the ball of radius $T^{-6}$ about $(Z,Z^*)$
satisfies these hypotheses with $P=2T\ge8$, so $\norm{\gamma_Z}\le(2T)^{12}\le T^{18}$ (as $2^{12}\le T^6$), and the two-variable matrix is holomorphic and bounded by $(2T)^{12}$ on that ball; Schwarz's lemma bounds its derivative at the centre by
$2(2T)^{12}T^6\le T^{25}$ (as $2^{13}\le T^7$). (d) $Y=Z(1+Z^{-1}(Y-Z))$ with $\norm{Z^{-1}(Y-Z)}\le T^{-5}\le1/2$ gives the inverse with $\norm{Y^{-1}}\le2T$, and
$Y^{-1}-Z^{-1}=Y^{-1}(Z-Y)Z^{-1}$ has norm at most $2T^{-4}$. Hence $\abs{\xi(Y)-\xi(Z)}\le2T^{-4}$, and since $\nu(Z)^2\ge T^{-2}$, $\nu(Y)^2\ge\nu(Z)^2-4T^{-4}\ge(1-8T^{-2})\nu(Z)^2$; the other bounds follow. For $\mathfrak g=\abs{\det}^2\nu^2$, write
$Z=(1-E)Y$ with $E=(Y-Z)Y^{-1}$, $\norm E\le2T^{-5}$; then $\abs{\det Z}\le(1+2T^{-5})^n\abs{\det Y}$, and
$\mathfrak g(Y)\ge(1+2T^{-5})^{-2n}(1-8T^{-2})\,\mathfrak g(Z)\ge e^{-4n/T^5}\cdot\frac78\,\mathfrak g(Z)\ge\frac12\mathfrak g(Z)$ for $T\ge8$ and $T\ge n$. (e) The operator norm of $Z$ is at most $nM^2\le M^3$.
By \eqref{eq:disc-nu}, with $\abs{\det Z}\le M^{3n}$ and $\nu(Z)^2\le1$, the inequality $M^{-1}\le\abs{\det Z}^2\nu^2$ bounds each factor below: $\abs{\det Z}^{-1}\le M^{1/2}$ and $\nu^{-2}\le M^{6n+1}$, so $\nu^{-1}\le M^{3n+1}$.
Cramer's rule gives $\norm{Z^{-1}}\le\abs{\det Z}^{-1}\norm{\mathrm{adj}\,Z}\le M^{1/2}\,n!\,M^{2(n-1)}\le M^{8n}$, as $n!\le M^n$. The second statement is (d) with $T=P/2$.
\end{proof}

\begin{theorem}[Mean value inequality]\label{thm:mean-value}
For measurable $f\ge0$ on $G_1$,
\[
  \int_\Omega c_G\,\mathfrak g(Z)^{k-1}\int_{U^0H}\int_{U^0H}f(s'\gamma_Zs)\,ds\,ds'\,dZ\le\int_{G_1}f ,
\]
with $c_G=\vol\U(k)/\vol\U(r)$ times the area of the unit sphere of $\C^k$. The law of $s\gamma_Z$ with density proportional to $\mathfrak g(Z)^{k-1}$ is at most the fibre measure of the
Gram map $Y\mapsto{}^t\bar YJ_1Y$ over the Gram matrix of the frame.
\end{theorem}
\lean{New.RealPlace.lintegral_gramSection_le, New.RealPlace.leGramFibre_gramSectionMeasure, New.RealPlace.lintegral_eq_lintegral_sphereMass_mul_dualFrameOfGram, New.RealPlace.dualFrameOfGram, New.RealPlace.gramSectionMeasure, New.RealPlace.sphereMass, New.Shape.openCellConst, New.Shape.openCellConst_pos, New.Shape.scaledVolumeRatio}

\begin{proof}
\emph{Dual tuples.} The tuples of $n$ vectors dual to the frame are $J_1F+Rp$, where $R$ has the columns $e_b$ and $x_i$ and $p$ runs over matrices. Every such tuple is $s\,y(T)$ for its Gram matrix $T$ and an element $s$ of the
fixer $U^0H$, where $y(T)$ is the standard dual tuple $e^*_b+\frac12\sum_aT_{ab}e_a$ ($b<n$), $-z+\sum_aT_{a,t}e_a+\sqrt{1+T_{tt}}\,x_{k_0}$ (defined when $1+T_{tt}\ge0$). Choosing first the $\varrho$ vectors $e^*_b+\dots$, which fixes the upper left block of $T$
and the unipotent part of $s$, and then the last vector, whose part in $X_0$ runs over a sphere of squared radius $1+T_{tt}$ (polar coordinates in $X_0$, with $H$ rotating it), Lebesgue measure is
\[
  dp=2^{-\varrho}\,m_k(1+T_{tt})\,dt\,ds,\qquad m_k(\rho)=\frac{\pi^k\rho^{k-1}}{(k-1)!}\ (\rho>0),
\]
with $t$ the real coordinates of $T$ and $ds$ Lebesgue measure on the unipotent parameters times the Haar probability measure of $H$.

\emph{Frames with given products.} For invertible $Z$, the tuples $y$ whose products with the frame form $Z$ are $y=y'Z$ with $y'$ dual to the frame, of Gram matrix $T'=Z^{-*}TZ^{-1}$ if $y$ has Gram matrix $T$. The substitution
costs $\abs{\det Z}^{2(\varrho+k)}$ in $dp$ and $\abs{\det Z}^{2(\varrho+1)}$ in $dT$, so $dp=\abs{\det Z}^{2(k-1)}2^{-\varrho}m_k(1+T'_{tt})\,dt\,ds$. At the Gram matrix $T_0=D$ of the frame, $T'=-\zeta^*\zeta$ and $1+T'_{tt}=\nu(Z)^2$, the dual tuple is
$F^\vee(Z)$, and $y(T')Z=\gamma_ZF$. By \eqref{eq:disc-nu}, $\abs{\det Z}^{2(k-1)}2^{-\varrho}m_k(\nu^2)=2^{-\varrho-1}\sigma_k\,\mathfrak g(Z)^{k-1}$, with $\sigma_k=2\pi^k/(k-1)!$ the area of the unit sphere of $\C^k$.

\emph{Below the fibre measure.} For a tuple $E$ of vectors, say that a measure $m$ on $G_1$ is \emph{below the fibre measure over $E$} if, for every lower semicontinuous $K\ge0$ on tuples $Y$ of the size of $E$ with
$\int_{T(Y)\in S}K\,dY\le C\vol S$ for all small open sets $S$ of hermitian matrices near the Gram matrix of $E$, one has $\int K(gE)\,dm(g)\le C$; here $T(Y)={}^t\bar YJ_1Y$ in its real coordinates and $dY$ is Lebesgue measure. This
is the sense in which the image of $m$ under $g\mapsto gE$ is at most the fibre measure $dY/dT$ of the Gram map \cite{Siegel1935}.

\emph{The section.} Let $\lambda$ be the law of $s\gamma_Z$ under $2^{-\varrho-1}\sigma_k\mathfrak g(Z)^{k-1}dZ\,ds$, $Z\in\Omega$. It is below the fibre measure over $F$: in the coordinates $(T,Z,s)$ above, the tuples with $Z\in\Omega$ have Lebesgue
measure with the density computed, continuous in $T$ at $T_0$, and the tuples with $Z\notin\Omega$ only add to $\int_{T(Y)\in S}K\,dY$.

\emph{From frames to Haar measure.} Three facts on the fibre measure give the comparison. (i) On bases, a measure below the fibre measure over the standard basis is at most $\vol\U(r)/2^r$ times Haar measure of $G_1$, normalised through its
compact dual; this is the computation of the fibre measure of the Gram map on bases. (ii) If $m$ is below the fibre measure over $E$, and $\nu$ fixes $E$ almost everywhere and maps, by $g\mapsto gE'$, to a measure below the fibre measure
of $Y'\mapsto({}^t\bar EJ_1Y',T(Y'))$ over $E'$, then $m*\nu$ is below the fibre measure over the joined tuple $(E,E')$, because Lebesgue measure on $(Y,Y')$ is the product. (iii) For $E'$ the dual vectors $e^*_b$, the measure
$2^{-\varrho}du$ on $U^0$ satisfies the hypothesis of (ii), by the coordinates of the dual tuples above; and for $E'$ the basis $x_i$ of $X_0$, the complement of the tuple $(f,e^*)$, which is definite, the measure $\vol\U(k)2^{-k}dh$ on $H$ does,
by (i) for $\U(k)$. Applying (ii) twice to $m=\lambda$ and then (i),
\[
  \lambda*2^{-\varrho}du*\frac{\vol\U(k)}{2^k}dh\le\frac{\vol\U(r)}{2^r}\,dg,\qquad\text{so}\qquad\lambda*ds\le\frac{\vol\U(r)}{2^{\varrho+1}\vol\U(k)}\,dg,
\]
as $r=k+2\varrho+1$. Integrating $f$ against $\lambda*ds$ gives the triple integral of the theorem with density $2^{-\varrho-1}\sigma_k\mathfrak g(Z)^{k-1}$; multiplying by $2^{\varrho+1}\vol\U(k)/\vol\U(r)$ turns it into $c_G\,\mathfrak g(Z)^{k-1}$ with
$c_G=\sigma_k\vol\U(k)/\vol\U(r)$. The second statement of the theorem is the paragraph on the section.
\end{proof}

\begin{theorem}[Open cell bound]\label{thm:open-cell}
\begin{enumerate}[label=\textup{(\alph*)}]
\item Let $F$ satisfy Schur's bound with constants $P\ge8r$ and $D$, let every $Y$ within $P^{-6}$ of $Z$ be bounded in the cell by $P$ with $\mathfrak g(Y)\ge c>0$, and let $\varphi$ be integrable,
vanishing where $\norm y>P^{20}$. Then $\int\norm{\int\varphi(y)F(\gamma_Z(uh)^{-1}y)dy}^2$ over $(h,u)\in H\times U^0$ with $\norm u\le P^{20}$, for Haar measure times $c_Gc^{k-1}$, is at most
$P^{2000n^2}D\int\abs\varphi^2$.
\item For a Bessel system (\cref{def:bessel-system}) $X$ of scale $B$, if $\mathfrak g(Z(\gamma))>0$ and some $y$ has $y$, $\gamma^{-1}y$ in the frame box of radius $\rho\ge B$, there are a compact $A$,
shifts in the fixer of the frame averaged by a probability measure, and a measure $\nu$ of mass at most
\[
  \Bigl(\frac{\rho^{9n}}{\mathfrak g}\Bigr)^{10^5n^4}\frac{(\#I)^2}{\deg\cdot c_G\cdot(\mathfrak g/2)^{k-1}}
\]
such that the mean of $\norm{\int_A\varphi(z)W(\gamma tz)\otimes\overline{W(tz)}dz}^2$ is at most $\int\abs\varphi^2d\nu$ for bounded measurable $\varphi$.
\item For systems $X_v$ at all infinite places and $\varphi$ of modulus at most $1$, invariant under the fixers, with $\varphi\cdot(W^\flat(y)\otimes W^\flat(\gamma^{-1}y)^*)$ supported in the boxes, the integral of
this product has compact support and norm at most the root of the product of the local bounds of (b), plus $\int\abs\varphi\cdot\norm{W\otimes W^*-W^\flat\otimes W^{\flat*}}$.
\end{enumerate}
\end{theorem}
\lean{New.RealPlace.lintegral_gramSection_integral_le, New.lintegral_mul_le_of_frameMatrix_le, New.RealPlace.HasDerivWords, New.RealPlace.HasDerivWords.exists_hasDerivAt_gramSection_mul, New.RealPlace.sobolev_const_mul_le_pow, New.RealPlace.lintegral_gramSection_le_of_hasDerivWords, New.RealPlace.convolutionWeight, New.RealPlace.hasDerivWords_integral, New.volume_upperEntries_closedBall_le, New.Shape.fixerMeasure, New.Shape.fixerBall, New.RealPlace.isBoundedInCell_of_inv_le_gramDiscReal, New.RealPlace.gramDiscReal_div_two_le_and_isBoundedInCell, New.BesselSystem.IsOpenCellBound, New.BesselSystem.exists_isOpenCellBound, New.BesselSystem.IsKernelBound, New.BesselSystem.exists_isKernelBound_gramSection, New.isL2Contraction_of_columns, New.BesselSystem.isL2Contraction_gramSection, New.pairingBoundSq, New.BesselSystem.pairKernel, New.BesselSystem.piKernel, New.BesselSystem.enorm_integral_pairKernel_cut_le, New.BesselSystem.exists_isCompact_enorm_integral_piKernel_le, New.BesselSystem.hasCompactSupport_and_enorm_integral_piKernel_le}

The proof of (a) rests on three lemmas.

\begin{lemma}[Unfolding along the frame]\label{lem:unfold-frame}
Let $w,\phi\ge0$ be measurable on $G_1$, $w$ left invariant under $R$ and $\phi$ of integral at most $K$ on every orbit of $U^0$, and let $R_0\ge1$ bound the entries of $m(s)$ for every $s\in U$ with $\phi(y)\ne0\ne\phi(sy)$ for some $y$. Then
$\int_{G_1}\phi w\le K(2R_0)^{2n^2}\int_{U\backslash G_1}w$. If $Q$ is a measurable set of matrices bounded in the cell by $P\ge1$ with $\mathfrak g\ge c>0$ on $Q$, then
\[
  \int_Q\int_{H\times U^0}w\bigl(\gamma_Z(uh)^{-1}\bigr)\,c_Gc^{k-1}\,du\,dh\,dZ\le(2P^2)^{2n^2}\int_{U\backslash G_1}w .
\]
\end{lemma}

\begin{proof}
Unfold along $U$; it suffices to bound $\int_U\phi(sy)\,ds$. Write $s=u\,\tilde x\,u_0$ with $u\in U^0$, $\tilde x$ the unitriangular lift of the entries $x$ above the diagonal of $m(s)$, and $\phi(u_0y)\ne0$: the integral over $u$ is at most $K$, and it
vanishes unless all coordinates of $x$ lie in the disc of radius $R_0$; the box of these coordinates has volume at most $(2R_0)^{2n^2}$. For the second statement, let $\chi$ be a cutoff function for $U^0$ and $\phi$ its product with the
indicator of $\{y:Z(y)\in Q\}$; compute $\int\phi w$ once in the coordinates of \cref{thm:mean-value}, where on $Q$ the density of the open cell is at least $c_Gc^{k-1}$ and along an orbit of the fixer only $\chi$ changes, with integral $1$; and once by
the first statement. If $Z(y)$ and $Z(sy)$ are bounded in the cell by $P$ and $s\in U$, the entries of $m(s)$ are at most $P^2$, since $Z(y)=m(s)^*Z(sy)$ by \cref{lem:frame-vectors}(b).
\end{proof}

\begin{lemma}[Sobolev at the section]\label{lem:sobolev-section}
Let $(V_w)$ be a family of measurable functions indexed by words in an alphabet $\mathcal L$, with norms left invariant under $R$, such that along a curve with $\gamma'=\gamma X$, $\norm X\le C$, the derivative of
$V_w\circ\gamma$ exists with norm at most $\Lambda C\sum_{\lambda\in\mathcal L}\norm{V_{\lambda w}\circ\gamma}$, where $\Lambda\ge0$ and $\#\mathcal L\,\Lambda\le P^{44}$. Let $P\ge8r$, let every $Y$ within $P^{-6}$ of $Z$ be bounded in the cell by $P$ with $\mathfrak g(Y)\ge c>0$, and
let $\int_{U\backslash G_1}\norm{V_w}^2\le K$ for $\abs w\le2n^2$. Then
\[
  \int_{\norm u\le P^{20}}\bigl\|V_\emptyset\bigl(\gamma_Z(uh)^{-1}\bigr)\bigr\|^2\,c_Gc^{k-1}\,du\,dh\le P^{550n^2}K.
\]
\end{lemma}

\begin{proof}
Choose real linear coordinates $Y=Z+e(x)$ on the matrices of size $n$, preserving Lebesgue measure, whose coordinate vectors have norm at most $P^2/4$, such that the cube $[0,P^{-8}]^{2n^2}$ lies within $P^{-6}$ of $Z$. Along a
coordinate line $t\mapsto\gamma_{p(t)}k$ with $\norm k\le P^{20}$, the logarithmic derivative of $\gamma_{p(t)}k$ has norm at most $P^{85}$ by \cref{prop:gram-section}(c),(d), so the derivatives of $t\mapsto V_w(\gamma_{p(t)}k)$ are bounded by
$\Lambda P^{85}$ times the next members of the family. Sobolev's inequality on the cube (the value at a corner bounded by the mean squares of the mixed derivatives of order at most one in each coordinate) costs at most
$2P^8(1+(\#\mathcal L\Lambda P^{85})^2)$ per coordinate; integrating over the cube with \cref{lem:unfold-frame} costs $2P^2$ per coordinate. The product over the $2n^2$ coordinates is at most $(P^{275})^{2n^2}=P^{550n^2}$.
\end{proof}

\begin{lemma}[Derivatives of convolutions]\label{lem:convolution-derivatives}
Let $V_w$, for words $w$ in the matrix units, be bounded continuous functions on $G_1$ such that the derivative of $V_w\circ\gamma$ along $\gamma'=\gamma X$ is $\sum_{a,b}X_{ab}V_{(a,b)w}\circ\gamma$. Let $\psi$ be integrable weights on $G_1$, indexed by words in
the letters $((a,b),c,d)$, with $\psi_{((a,b),c,d)i}(y)=\psi_i(y)\,(y^{-1})_{ad}\,y_{cb}/R_0^2$. Then the functions $g\mapsto\int\psi_i(y)V_{i'}(gy)\,dy$, with $i'$ the word of the first components of $i$, satisfy the hypothesis of
\cref{lem:sobolev-section} with $\Lambda=R_0^2$.
\end{lemma}

\begin{proof}
Let $G_i(g)=\int\psi_i(y)V_{i'}(gy)\,dy$ and let $\gamma$ be a curve with $\gamma'=\gamma X$, $\norm X\le C$. For fixed $y$ the curve $\gamma(t)y$ has $(\gamma y)'=(\gamma y)(y^{-1}Xy)$, so by hypothesis
\[
  \frac d{dt}V_{i'}(\gamma(t)y)=\sum_{a,b}(y^{-1}Xy)_{ab}\,V_{(a,b)i'}(\gamma(t)y),\qquad (y^{-1}Xy)_{ab}=\sum_{c,d}(y^{-1})_{ad}X_{dc}y_{cb}.
\]
Multiplying by $\psi_i(y)$ and writing $\lambda=((a,b),c,d)$, the factor $\psi_i(y)(y^{-1})_{ad}y_{cb}$ is $R_0^2\,\psi_{\lambda i}(y)$, and the first components of $\lambda i$ form the word $(a,b)i'$. The derivatives are dominated by
$R_0^2C\sum_\lambda\abs{\psi_{\lambda i}}\sup\norm{V_{(a,b)i'}}$, which is integrable, so we may differentiate under the integral:
\[
  \frac d{dt}G_i(\gamma(t))=\sum_\lambda R_0^2X_{dc}\,G_{\lambda i}(\gamma(t)),
\]
whose norm is at most $R_0^2C\sum_\lambda\norm{G_{\lambda i}(\gamma(t))}$, because $\abs{X_{dc}}\le\norm X\le C$.
\end{proof}

\begin{proof}[Proof of \cref{thm:open-cell}]
(a) Let $V_w$ be the right derivatives of $F$ given by \cref{thm:bessel-schur}(a), and $\psi_i$ the weights of \cref{lem:convolution-derivatives} built from $\varphi$ with $R_0=P^{20}$; as $\varphi$ vanishes where $\norm y>P^{20}$, $\abs{\psi_i}\le\abs\varphi$. By Schur's
bound the convolutions $g\mapsto\int\psi_i(y)V_{i'}(gy)\,dy$ have squared norms at most $P^{2\abs i}D\int\abs\varphi^2\le P^{4n^2}D\int\abs\varphi^2$ for $\abs i\le2n^2$. There are $r^4$ letters and $\Lambda=P^{40}$, so $\#\mathcal L\Lambda\le P^{44}$, and
\cref{lem:sobolev-section} gives the bound $P^{550n^2}P^{4n^2}D\int\abs\varphi^2\le P^{2000n^2}D\int\abs\varphi^2$.

(b) Put $Z=Z(\gamma)\in\Omega$ and $M=\rho^{9n}/\mathfrak g(Z)$. The entries of $Z$ are at most $\rho^2$ by \cref{lem:frame-vectors}(b), so $\mathfrak g(Z)\le\rho^{8n}$ and $\rho\le M$, $M^{-1}\le\mathfrak g(Z)$; let $M_0=M^{9n}$ and $P=(M_0^5)^n$. By \cref{prop:gram-section}(b),
$\gamma t(h,u)=s\gamma_Z(uh)^{-1}$ for some $s\in R$ and a measurable $t\colon H\times U^0\to U^0H$. Take $A$ a compact tube over a set $\Omega_0\subseteq U^0$ of elements of norm at most $P^{20}$, the probability measure $m$ the fixer
measure conditioned on $\Omega_0$, and $\nu$ the measure $\norm W^2\,dz$ on $A$ times $(m(\Omega_0))^{-1}P^{2000n^2}\#I/\deg$. The Bessel system satisfies Schur's bound with $P\ge B^5$ and $D=\#I/\deg$ (\cref{def:bessel-system}); the near-point bounds
of \cref{prop:gram-section}(e) give the hypothesis of (a) with $c=\mathfrak g(Z)/2$. For each column of the kernel, (a) bounds the mean over $m$ of $\norm{\int_A\varphi(z)W(s\gamma_Z(uh)^{-1}z)\overline{W_j(tz)}\,dz}^2$; the norm of
$\int\psi(y)W(sgy)\otimes\overline{W(ty)}$ does not depend on $s,t\in R$. The mass of $\nu$ is at most $M^{10^5n^4}(\#I)^2/(\deg c_G(\mathfrak g/2)^{k-1})$, by $\int_A\norm W^2\le K(2R^2)^{2n^2}\#I$ for the measure $K$ of the fibres of $A$ along $U^0$, and
the identity $P^{2000n^2}M_0^5(MP^{40})^{2n^2}=M^{93600n^4+2n^2+45n}\le M^{10^5n^4}$. Finally, if $z$ is in the frame box of radius $\rho\le M$ and $\ell(s\gamma_Z(uh)^{-1}z)\le M$, then $\norm{u^{-1}z}\le M_0^{5}$: move $\gamma_Zz$ by an element of $U$ to a point of
norm at most $M$, and express the columns $z^{-1}e^*_a$ through the frames at both points with coefficients from $Z^{-1}$; so $z\in A$ for the tube of radius $M_0^5$.

(c) At each place (b) gives a compact $A_v$, measures $\nu_v$ and probability measures $m_v$ on the shifts. The operator bounds multiply over the places, so the mean over the product of the $m_v$ of the squared norm of the
pairing is at most $\int\abs\varphi^2d\prod\nu_v\le\prod_v\nu_v(G_1)$, as $\abs\varphi\le1$; hence some shift $t$ does no worse than the mean, and $\varphi$, being invariant under the fixers, is unchanged by it. After the translation by $t$ the cut
integrand lives on the compact set $\prod A_v$, where the difference with the uncut integrand is bounded by $\int\abs\varphi\,\norm{W\otimes W^*-W^\flat\otimes W^{\flat*}}$. The kernel $W(y)\otimes W(\gamma^{-1}y)^*$ is the adjoint of $W(\gamma^{-1}y)\otimes W(y)^*$, and
$\gamma^{-1}$ has the same discriminant as $\gamma$.
\end{proof}

\subsection{Bessel systems}\label{ss:bessel-systems}

\begin{definition}[Bessel system]\label{def:bessel-system}
A \emph{Bessel system} of scale $b\ge8r$ and decay rate $f$ consists of a finite set $I$, continuous functions $W$, $W^\flat\colon G_1\to\C^I$, a degree $\deg>0$ and an irreducible
unitary representation $\pi$ of $H$ on $\C^I$, such that $(W,W^\flat)$ is a cut pair of type $\pi$ with decay rate $f$ (\cref{thm:bessel-schur}(c)), $W$ satisfies Schur's bound with
constants $P$ and $\#I/\deg$ for every $P\ge b^5$, and $W^\flat$ has the \emph{local bounds} of \cref{thm:local-band-limits}: a sieve system of forms whose top form is the inner product
of normalised orbital integrals of $W^\flat$.
\end{definition}
\lean{New.BesselSystem, New.HasLocalBounds, New.FullOrCut.select, New.besselOrbitalIntegral, New.normalizedOrbitalIntegral, New.RealPlace.gramDiscReal, New.Shape.openCellConst}

\begin{theorem}[Bessel systems exist]\label{thm:bessel-systems-exist}
For admissible data and each infinite place $v$ there is a Bessel system of scale $b$ and decay rate $f$ whose full function is the Bessel function of the discrete series of
parameter $2\varsigma(v)$ (the target parameter, doubled) and whose type has the column numbers of the weight at $v$. Its degree is the formal degree, so that
$\deg_v$ is the dimension of the corresponding representation of the compact form, and $\log\prod_v\deg_v$ lies within explicit bounds; each system has at most $b^{2k}$
coordinates. The filling data of $D$ consist of $N$, $p_0$, $p_1$ with their chosen places, $\pi_N$, the generator $\varpi$ of the different, the level $K_0$ of $X_0$ and these
systems; the level $K_0$ is the preimage of the principal congruence subgroup $K(M)$, which includes the congruence modulo $\pi_N$.
\end{theorem}
\lean{New.Data.exists_besselSystem, New.Data.filling, New.Data.isDiscreteSeriesBessel_filling_besselSystem, New.Data.levelX_zero_eq_comap, New.Data.card_X₀, New.Data.unitarySignFrameY, New.Data.exists_unitary_schur_rep, New.Data.weylDimensionPoly_paramW_le, New.Data.volU_div_volU_le, New.Data.scaledVolumeRatio_le, New.Data.weylDimensionPoly_mul_scaledVolumeRatio_le, New.Data.abs_paramW_le, New.Data.exists_lowestKType_eq_ite, New.Data.realPlaceSizes, New.Data.exponent_add_le_mul_sq_r, New.Data.liftWeight_add_columnCount_le, New.Data.rank_add_le_cube, New.Data.rank_sq_mul_log_le, New.IsBesselFunction.integral_dsKernel_mul_eq_of_convolve_eq, New.Data.natCast_levelPrime_dvd_iff}

\begin{proof}
\emph{The type.} Let $\pi$ be the representation of $H=\U(X_0)$ given by the Schur matrices of the weight with the column numbers of the weight at $v$, written in an orthonormal basis for its invariant Gram form (write the positive
Gram matrix as $B^*B$ and conjugate by $B$). It is continuous, unitary and irreducible, and at a highest weight unit vector $e$ its coefficient $\langle e,\pi(h)e\rangle$ is the discrete series coefficient $c_u$ of the
parameter $u$ of this weight, a product of principal minors in the decreasing order of $u$. The lowest weight of $u$ is the number of columns of height one at a single index and $0$ elsewhere, and the entries of $u$ are at most $3a$.

\emph{The parameter.} Let $l=2\varsigma(v)$ in a unitary frame with the signs of $Y_0$. Its entries are at most $3a$; it is steep for $B$ (\cref{prop:grids}) with $3B$ plus the column number at most $b^2$; its gap exceeds $500r^3$ by at least
$r^7$, and $r^7\ge10^{30}((r+k)^2\log b^3)^3$ by \cref{prop:sizes}. Its formal degree is Weyl's dimension polynomial at $l$, at most $b^{2r^2}$ since each of the at most $r^2$ halved differences is at most $3a\le b^2$.

\emph{The point.} With $\beta=b^3\ge r+k+3a$, \cref{thm:bessel-average}(c) gives $x$ with $\ell(x)\le b^{5n}$ at which $V(x)\le\beta^{20(r+k)^2}\abs{\text{period}}$, and by \cref{thm:bessel-average}(b) $V(x)\le b^{60(r+k)^2}\deg\norm{W_x}^2$. As $k<r$,
$m=60(r+k)^2\le240r^2\le301r^2$.

\emph{The functions.} \Cref{thm:bessel-schur}(b),(c) give $F=cW_x$, reproduced by $e$, and its cut, forming a cut pair; the exponents satisfy $3\cdot20(r+k)^2+99r^2\le400r^2$. Schur's bound for $F$ (\cref{thm:bessel-schur}(a)) holds with every
$P\ge b^5\ge(r+1)^4(3a+r+1+2n^2)$ and $D=\norm F_2^2/\deg=\#I/\deg$. The local bounds of the cut follow from \cref{thm:local-band-limits}, whose size conditions hold with box radius $b^{6n+2}$ and exponent $k-(3n+6)$; they use the norm of the
matrix coefficient of $F$ and the volume of balls. The number of coordinates is the dimension of $\pi$, at most $b^{2k}$, and $\vol\U(k)/\vol\U(r)\le b^{r^2}$, so $\deg_v$ times the constant of the local bounds is at most $b^{4r^2}$.

\emph{The data.} The filling data record $N$, the primes $p_0,p_1$ with the places dividing $\pi_M$, the idele of $\pi_N$, the generator of the different and these systems. The level $D.\mathrm{level}_0$ of $X_0$ is the set of $k$ that satisfy the
congruence modulo $\pi_N$ and reduce to $1$ at the chosen places; as $p_0\ne p_1$ are unramified, $p_0p_1$ divides an integral adele exactly when both primes do, which happens exactly when its components at the chosen places and their
conjugates lie in the maximal ideals. Hence the level is the preimage of $K(M)$, $M=p_0p_1$.
\end{proof}
